%====================================================
% 05 CNN Architecture Design
%====================================================

\section{CNN Architecture Design}

Convolutional Neural Networks (CNNs) are a class of deep learning models specifically designed for image processing and computer vision tasks. Unlike traditional artificial neural networks, CNNs automatically learn hierarchical image features through convolution and pooling operations, eliminating the need for manual feature extraction\cite{goodfellow2016deep,lecun1998gradient}.

In this assignment, a baseline CNN architecture was first developed for handwritten digit recognition using the MNIST dataset. Four additional CNN models were then created by modifying selected architectural parameters, including the convolution kernel size, the number of convolution layers, the optimization algorithm, and the activation function. All models were implemented using TensorFlow and Keras.

\subsection{Baseline CNN Architecture}

The baseline CNN model consists of the following layers:

\begin{enumerate}
    \item First Convolution Layer
    \item First Max Pooling Layer
    \item Second Convolution Layer
    \item Second Max Pooling Layer
    \item Flatten Layer
    \item Fully Connected Dense Layer
    \item Output Layer
\end{enumerate}

The architecture receives a grayscale image of size $28 \times 28 \times 1$ as input and produces the probability of each handwritten digit class (0--9).

\subsection{Layer Description}

\subsubsection{Convolution Layer}

The convolution layer is responsible for extracting important image features such as edges, curves, and local patterns. Each convolution filter slides across the input image and generates a feature map.

The convolution operation is defined as

\begin{equation}
Y(i,j)=\sum_{m}\sum_{n}X(i+m,j+n)K(m,n)
\end{equation}

where

\begin{itemize}
    \item $X$ is the input image,
    \item $K$ is the convolution kernel,
    \item $Y$ is the extracted feature map.
\end{itemize}

The baseline model uses 32 filters in the first convolution layer and 64 filters in the second convolution layer.

\subsubsection{Max Pooling Layer}

Max pooling reduces the spatial dimensions of the feature maps while preserving the most significant information. A pooling window of $2\times2$ with a stride of 2 was used throughout the models.

This operation decreases computational complexity and reduces the likelihood of overfitting.

\subsubsection{Flatten Layer}

After the convolution and pooling operations, the extracted feature maps are transformed into a one-dimensional feature vector using a Flatten layer.

This vector becomes the input to the fully connected neural network.

\subsubsection{Fully Connected Layer}

The fully connected layer performs high-level reasoning using the extracted image features.

The baseline architecture contains

\begin{itemize}
    \item 128 neurons
    \item ReLU activation function
\end{itemize}

The number of neurons remained unchanged throughout all five models to ensure a fair comparison.

\subsubsection{Output Layer}

The output layer consists of ten neurons corresponding to the ten digit classes (0--9).

A Softmax activation function is applied to convert the output values into probabilities.

The predicted class is obtained by selecting the neuron with the highest probability.

\subsection{Model Compilation}

All CNN models were compiled using the Sparse Categorical Crossentropy loss function because handwritten digit recognition is a multi-class classification problem.

Model performance was evaluated using classification accuracy.

The baseline compilation configuration is summarized below.

\begin{itemize}
    \item Optimizer: Adam
    \item Loss Function: Sparse Categorical Crossentropy
    \item Evaluation Metric: Accuracy
\end{itemize}

The optimizer and activation function were modified only in the respective comparison models discussed in later chapters.

\subsection{Baseline CNN Architecture Summary}

Figure~\ref{fig:model_summary} presents the summary of the baseline CNN architecture generated by TensorFlow Keras.

\begin{figure}[H]
\centering
\includegraphics[width=0.90\textwidth]{step_2_2_model_summary.png}
\caption{Baseline CNN model summary generated using TensorFlow Keras.}\label{fig:model_summary}
\end{figure}

The baseline CNN contains two convolution layers, two max pooling layers, one flatten layer, one fully connected dense layer with 128 neurons, and an output layer with ten neurons. The total number of trainable parameters is 225,034.

\subsection{Architecture Variations}

To investigate the influence of different CNN design parameters, four additional models were developed based on the baseline architecture. Each model modifies only one parameter while keeping all other settings unchanged.

\begin{table}[H]
\centering
\caption{Summary of CNN Models Developed in this Assignment}\label{tab:cnn_models}

\begin{tabular}{llll}
\toprule
\textbf{Model} & \textbf{Modified Parameter} & \textbf{Configuration} & \textbf{Purpose}\\
\midrule
Model 1 & Baseline & Standard CNN & Reference model\\
Model 2 & Kernel Size & $5\times5$ kernels & Kernel comparison\\
Model 3 & CNN Depth & Three convolution layers & Feature extraction comparison\\
Model 4 & Optimizer & SGD optimizer & Optimization comparison\\
Model 5 & Activation Function & Tanh activation & Activation comparison\\
\bottomrule
\end{tabular}

\end{table}

These five CNN architectures were trained under identical experimental conditions and compared using testing accuracy, testing loss, training time, confusion matrix analysis, misclassification analysis, and custom handwritten digit prediction. The detailed experimental setup and comparative analysis are presented in the subsequent chapters.